\section{The acquisition--authority envelope}
\label{sec:envelope}

The acquisition/closure split becomes substantially more useful when it is
lifted from a design rule to a comparison theory.  A scientific discovery
controller has four coupled responsibilities: generate or select a route,
acquire evidence through it, convert acquired material into decision-usable
state while preserving useful future cues, and decide whether the task may
close.  A retriever improves only one part of this object; a stopping policy can
save cost while making another part worse.  This section defines the joint
object and derives the conditions under which such local improvements can be
composed into a scientifically stronger controller.

\paragraph{Mathematical conventions.}
Unless a probability law is explicitly introduced, exactness and
factorization are pointwise, set-theoretic statements. Fibre constancy alone
therefore proves a set map, not a measurable or computable implementation; a
measurable factorization additionally requires a measurable quotient/factor
map. Probabilistic statements are almost-sure under their declared laws and do
not certify null fibres. Decision and action alphabets are nonempty. When total
variation is used below, both laws live on one declared measurable transcript
space and
$\operatorname{TV}(P,Q)=\sup_{B}|P(B)-Q(B)|$, one half of the $L^1$ distance
when densities exist.

\subsection{Worlds, obligations and a non-compensatory frontier}

Let $w\in\mathcal W$ be a scientific world and let $h_t$ be the complete
observable transcript after charged step $t$.  The task contract assigns to
$w$ a set $\mathcal O(w)$ of \emph{material scientific obligations}.  An
obligation may be a claim that needs evidence, a contradiction that needs both
sides inspected, a dataset or code artifact needed to establish reproducibility,
a route whose availability must be resolved, or a previously retrieved source
that must be processed under the current question.  Materiality means that
resolving the obligation can change the predeclared decision, conclusion,
uncertainty state or admissible closure state.  It is not synonymous with
relevance to a search query.

A controller $\pi$ comprises route generation $G$, acquisition $A$, a
question-conditioned state compiler $M$, and closure authority $C$.  Under
budget $B$, source-access contract $\mathcal I$ and task distribution $P$, let
its outcome vector be (with the admissible policy class restricted to policies
for which these coordinates are measurable and integrable)
\[
 Z_B(\pi,w)=(D,U,V,-F,Y,-K).
\]
Here $D$ is decision-relevant evidence discovery, $U$ is final scientific
utility, $V$ is the value of search options preserved for prospectively hidden
future questions, $F$ is false task closure, $Y$ is valid closure yield, and
$K$ is fully charged cost.  Signs are chosen so larger is better except that
$F$ is a non-compensatory safety coordinate: a gain in report quality cannot
purchase permission to hide a false closure.

For two outcome vectors $z,z'\in\mathbb R^6$, write $z\succeq z'$ for the
coordinatewise product order,
\[
 D\geq D',\quad U\geq U',\quad V\geq V',\quad
 F\leq F',\quad Y\geq Y',\quad K\leq K',
\]
and write $z\succ z'$ when $z\succeq z'$ and at least one inequality is
strict. Thus a utility gain never compensates for worse false closure or cost;
incomparability is retained as a scientific trade-off rather than collapsed by
an unregistered scalar score. Let $q(\pi)$ collect the nonnumeric provider,
observability, custody and contract-validity states of policy $\pi$. For a
prospectively fixed guard registry
$\eta=(\tau_F,\tau_Y,B,q_{\rm req})$, define
\[
 \mathsf G_\eta(\pi)=1
 \quad\Longleftrightarrow\quad
 \mathbb E[F]\leq\tau_F,
 \ \mathbb E[Y]\geq\tau_Y,
 \ \mathbb E[K]\leq B,
 \ \text{and }q(\pi)=q_{\rm req}.
\]
An unbound or \textsc{cannot-check} guard is not a pass.

\begin{definition}[Acquisition--authority envelope]
\label{def:acquisition-authority-envelope}
For an admissible policy class $\Pi$,
\[
 \mathcal E_{B,\mathcal I}(\Pi)
 =\operatorname{cl}\left\{
     \mathbb E_{w\sim P}[Z_B(\pi,w)] :
     \pi\in\Pi,\ \mathbb E[K(\pi,w)]\le B
   \right\}.
\]
Its \emph{guarded frontier} retains only points satisfying the prospectively
fixed guard predicate $\mathsf G_\eta$ and then takes the $\succeq$-maximal
points. Guard filtering and maximization are different operations: guarded
attainable sets are monotone under donor addition, whereas their maximal-point
sets need not be.
\end{definition}

\begin{definition}[Prospectively registered strict improvement]
\label{def:strict-envelope-improvement}
Before outcomes, register a donor-complete class $\Pi_0$, the contracts
$(P,B,\mathcal I)$, guard registry $\eta$, primary coordinate $j$, strict
margin $\delta_j>0$, coordinatewise non-inferiority tolerances
$\varepsilon\geq0$, and a simultaneous one-sided uncertainty rule. Let
$\mathcal F_0$ be the resulting guarded donor frontier. A candidate point
$z_c$ establishes \emph{strict donor-complete improvement} only if its guards
pass and, for every $z\in\mathcal F_0$,
\[
 z_{c,k}\geq z_k-\varepsilon_k\quad\text{for every benefit-signed coordinate }k,
 \qquad z_{c,j}\geq z_j+\delta_j,
\]
with the sign convention above and with the registered simultaneous bounds,
not point estimates, satisfying the inequalities. Safety coordinates may set
no-harm tolerance to zero and always remain subject to $\mathsf G_\eta$.
Failure to bind the donor frontier, any guard, or the uncertainty rule is
\textsc{cannot-check}. A new nondominated trade-off may be reported as frontier
extension, but neither that fact nor a win over a post-selected weak comparator
licenses the word ``superior.''
\end{definition}

This definition absorbs the strongest neighbouring components instead of
opposing them.  Lexical and corpus-side augmentation, reasoning-aware retrieval,
fielded inspection and selective fetching, and citation expansion enlarge the
action family of $G$ and $A$ \cite{sage2026,agentir2026,sieve2026,rethinkinglitsearch2026}.
Claim--evidence coverage, marginal information utility, conformal continuation,
decision-aware and evidence-aware stopping enlarge the policy and signal class
\cite{halt2026,deepcontrol2026,micp2026,decisionstop2026,donotstopearly2026,confidencebasedstop2026}.
Finite source certificates, obligation graphs and question-conditioned memory
enlarge $M$ and the observable closure state
\cite{knowplan2026,icore2026,memchain2026,scienceintent2026}.  These are donor
territory.  The residual question is whether their composition reaches a better
guarded frontier and whether its closure decision is identifiable from the
state the product actually retains.

\begin{proposition}[Donor-monotonicity]
\label{prop:donor-monotonicity}
If $\Pi_1\subseteq\Pi_2$ under identical task, access and budget contracts, then
$\mathcal E_{B,\mathcal I}(\Pi_1)\subseteq
\mathcal E_{B,\mathcal I}(\Pi_2)$.  Consequently, adding a sound donor action,
signal or memory operation cannot make the attainable envelope smaller, though
it may leave the guarded frontier unchanged after its cost is charged.
\end{proposition}

\begin{proof}
Every policy feasible in $\Pi_1$ remains feasible in $\Pi_2$; take the closure
of their attainable outcome vectors. Intersecting both attainable sets with the
same prospectively fixed guard predicate preserves inclusion. Taking maximal
points afterward need not preserve inclusion, which is why the proposition is
about attainable sets rather than frontier-point identity.
\end{proof}

The inclusion is an inclusion of closed attainable sets, not an inclusion of
their Pareto-maximal or guarded-frontier points. A newly attainable point can
dominate and thereby remove an old frontier point.

This elementary monotonicity result changes the novelty test.  A new paper that
improves retrieval should be absorbed into the comparator class.  The theory is
falsified, not protected, if the expanded donor class reaches every proposed
candidate point under matched resources. Any strict-superiority statement must
apply Definition~\ref{def:strict-envelope-improvement}; merely lying outside a
plotted frontier, beating one convenient member, or choosing the improved
coordinate after seeing outcomes is insufficient.

\subsection{The acquisition ceiling: why seven emitted identifiers matter}

Let $G_w$ be a complete evaluator gold set for a fixed task world, $A_\pi(w)$
the union of gold-normalized identities ever emitted by any admissible route
during a frozen run, and $S_\pi(w)$ the final submitted set.  A downstream
ranker, reader or synthesizer that is forbidden to invent evidence satisfies
$S_\pi(w)\cap G_w\subseteq A_\pi(w)\cap G_w$.

\begin{theorem}[Run-conditional acquisition ceiling]
\label{thm:acquisition-ceiling}
For any such downstream policy and any nonempty $G_w$,
\[
 \operatorname{Recall}(S_\pi;G_w)
 \leq \frac{|A_\pi(w)\cap G_w|}{|G_w|}.
\]
More generally, if obligations have nonnegative weights $\mu_o$ and each
obligation can be discharged only by at least one member of evidence set $E_o$,
then the unnormalized resolved-obligation weight is at most the total weight of
obligations for which $A_\pi(w)\cap E_o\ne\varnothing$.
\end{theorem}

\begin{proof}
The submitted relevant set is a subset of the acquired relevant set, so its
cardinality cannot be larger.  For the weighted form, an obligation whose
required evidence set is disjoint from the acquired pool cannot be discharged
by a downstream operation over that pool; summing the remaining nonnegative
weights gives the bound.
\end{proof}

The theorem separates candidate-generation failures from selection failures
without assuming that a particular retriever is optimal.  In the archived
external adverse campaign, all routes emitted only 7 of the 229 gold
identifiers; six were submitted and one was discarded.  The run-conditional
recall ceiling was therefore $7/229=0.030568$, and 222 gold identities were
unavailable to every downstream selection rule on that capture.  This does not
estimate general web discoverability.  It proves the narrower but decisive
causal fact that reranking or synthesis could not repair that run.  The adverse
terminal and the separate provider-validity failure remain unchanged.

The same analysis exposes why matched request counts are not a matched
scientific opportunity.  Let $\xi(a,G_w)$ be the admissible exposure of action
$a$: the prospectively validated opportunity, through the scorer-compatible
interface, to emit identities in $G_w$.  Equal $\sum_a 1$ does not imply equal
$\sum_a\xi(a,G_w)$.

\begin{proposition}[Call matching does not identify exposure matching]
\label{prop:call-exposure}
For every positive call budget $m$, there exist two $m$-call policies under the
same nominal budget whose run-conditional acquisition ceilings are respectively
one and zero.  Hence call-count equality alone gives no nontrivial bound on the
difference in discoverability.
\end{proposition}

\begin{proof}
Take a one-item gold set.  Give the first policy $m$ calls whose admissible
support contains the item with a deterministic first-call return, and give the
second $m$ calls to an interface whose output support is disjoint from the gold
identity namespace.  Both spend $m$ calls; their acquired-gold fractions are one
and zero.
\end{proof}

The archived campaign gave the lexical arm three calls to
the only scoring-admissible backend and the governed arm one such call plus two
inadmissible calls.  Its three-versus-three request count was therefore not an
exposure match.  In a future protected study, exposure must be certified
outcome-blind through interface-class witnesses or treated as
\textsc{cannot-check}; gold itself cannot be used as a worker-visible route
oracle.

\subsection{When can a product implement scientific closure?}

Let $s:\mathcal W\rightarrow\mathcal S$ be everything a donor product makes
available to its closure rule at a fixed decision time: acquired evidence,
coverage and utility signals, route states, compiled memory and charged
resource state.  Write
\[
 \mathcal Y_C=\{\textsc{close},\textsc{open},\textsc{cannot-check}\}
\]
and let $\kappa:\mathcal W\rightarrow\mathcal Y_C$ be the correct task
decision under the declared scientific contract.

\begin{theorem}[Closure factorization]
\label{thm:factorization}
There exists a decision map $g:\mathcal S\rightarrow\mathcal Y_C$ such that
$\kappa=g\circ s$ if and only if $\kappa$ is constant on every fibre
$s^{-1}(z)$.
\end{theorem}

\begin{proof}
If $\kappa=g\circ s$ and $s(w)=s(w')$, then
$\kappa(w)=g(s(w))=g(s(w'))=\kappa(w')$.  Conversely, if $\kappa$ is constant
on every fibre, first define $g(z)=\kappa(w)$ for $z\in s(\mathcal W)$ using
any $w$ with $s(w)=z$. Fibre constancy makes this definition independent of
the chosen representative. Then extend $g$ from $s(\mathcal W)$ to all of
$\mathcal S$ by assigning any fixed member of the nonempty alphabet
$\mathcal Y_C$ to every $z\notin s(\mathcal W)$. The extension does not affect
$g\circ s$.
\end{proof}

The theorem gives an exact necessary-and-sufficient condition, not a claim that
a particular architecture is intrinsically more expressive.  Coverage,
confidence and route-local stop signals are sufficient for task closure only
when no two worlds with the same retained state require different task
decisions.  A state that drops the materiality or unresolved status of an
unavailable route creates a mixed fibre and cannot be repaired by a different
classifier over that same state.  An information-equivalent product that
retains the missing coordinate has pure fibres and can tie.  This is exactly the
boundary exhibited by the protected exact-contract battery: the available-route
product fails on unresolved material routes, while the ideal typed product is
decision-identical to the authority rule.

\subsection{The price of an observationally unresolved world}

The factorization obstruction has a quantitative version.  Consider two task
worlds $w_0,w_1$ whose correct decisions are respectively \textsc{close} and
\textsc{open}. Let $P_0$ and $P_1$ be probability measures on a common
measurable transcript space $(\mathcal X,\Sigma_{\mathcal X})$, representing the
complete observable transcript available to the controller before its decision. A
randomized closure rule is a Markov kernel from
$(\mathcal X,\Sigma_{\mathcal X})$ to
the finite discrete alphabet $\mathcal Y_C$. Count
\textsc{cannot-check} as not yielding a valid closure in $w_0$ but not as false
closure in $w_1$.

\begin{theorem}[Indistinguishable-world lower bound]
\label{thm:indistinguishable}
For every randomized closure rule, its false-closure probability $F_1$ in
$w_1$ and valid-closure shortfall $L_0$ in $w_0$ satisfy
\[
 F_1+L_0\geq 1-\operatorname{TV}(P_0,P_1).
\]
In particular, if the two worlds are observationally identical, then
$F_1+L_0=1$ and $\max(F_1,L_0)\geq 1/2$.
\end{theorem}

\begin{proof}
The decision rule is a binary test between $P_0$ and $P_1$ once
\textsc{open} and \textsc{cannot-check} are combined as ``do not close.''  The
sum of its two testing errors is bounded below by one minus total variation.
When $P_0=P_1$, a transcript-dependent close probability has the same
expectation in both worlds, so false closure plus failure to close the closable
world equals one.
\end{proof}

This lower bound prevents a cosmetic conversion of every negative result into a
positive terminal.  If the decisive route is absent from the observable action
support, no stopping heuristic can create the missing distinction.  Scientific
progress must instead do at least one of three things: invent an acquisition
action that makes $P_0$ and $P_1$ distinguishable, retain a state coordinate
that an earlier compiler erased, or introduce and validate a substantive world
assumption that rules out one member of the pair.  Otherwise a calibrated open
or \textsc{cannot-check} state is the only honest output.

\subsection{Authority completion and its maximality}

For a transcript $h$, let $\mathcal W(h)$ be the nonempty class of task worlds
compatible with the declared observation and world contract, and let
$A(w,h)$ say that closing at $h$ is scientifically admissible in world $w$.
The robustly safe histories are
\[
 \mathcal H_{\rm safe}
 =\{h:\text{$A(w,h)$ holds for every $w\in\mathcal W(h)$}\}.
\]

\begin{theorem}[Maximally permissive robust completion]
\label{thm:maximal-completion}
The rule $C^\star$ that closes exactly on $\mathcal H_{\rm safe}$ has no false
closure on any compatible world. Every other rule with that universal soundness
property closes on a subset of $\mathcal H_{\rm safe}$. Hence $C^\star$ is the
unique maximally permissive universally sound close set, relative to the
declared compatible-world class.
\end{theorem}

\begin{proof}
Closure by $C^\star$ is admissible in every compatible world by the definition
of $\mathcal H_{\rm safe}$. If another universally sound rule closes at
$h\notin\mathcal H_{\rm safe}$, some $w\in\mathcal W(h)$ has $\neg A(w,h)$,
contradicting soundness. Thus every sound close set is a subset of
$\mathcal H_{\rm safe}$. \hfill$\square$
\end{proof}

The theorem is exact but not directly computable unless the compatible-world
class is represented adequately. The registered obligation rule is a
conservative implementation candidate. Let $R(h)$ be the set of material
obligations not yet independently discharged and let $V(h)$ say that every
required task-contract, observability and provider-validity condition is
established. Define
\[
 C^+(h)=
 \begin{cases}
  \textsc{cannot-check}, & \neg V(h),\\
  \textsc{close}, & V(h)\ \text{and}\ R(h)=\varnothing,\\
  \textsc{open}, & \text{otherwise}.
 \end{cases}
\]
It is sound whenever
\[
 V(h)\ \wedge\ R(h)=\varnothing\quad\Longrightarrow\quad
 h\in\mathcal H_{\rm safe}.
\]
It equals the maximal rule $C^\star$ if and only if the reverse implication
also holds. Equivalently, the registry and validity predicate must be complete
for robust closure, not merely sufficient. An adversarial-realizability
condition---every unresolved material obligation has a compatible world in
which closure is inadmissible---establishes that $R(h)\ne\varnothing$ implies
$h\notin\mathcal H_{\rm safe}$; it does not by itself prove that every safe
history is recognized by $V$ and $R$.

This distinction is the open-world boundary. A closed registry, a validated
sampling model, independently curated hidden-route custody, or a declared
\textsc{cannot-check} state may justify the missing equivalence. Without it,
$C^+$ is a sound lower approximation that may forgo valid closure, not a
maximally permissive rule. Future acquisition, ranking, stopping or memory
donors can be installed below $C^+$ without losing soundness when the sufficient
condition is preserved; improvements that establish $V$, discharge obligations,
or split a previously mixed closure fibre can move $C^+$ toward $C^\star$.

\begin{corollary}[When acquisition enlarges closure authority]
Adding a donor mechanism can expand discovery or reduce cost by
Proposition~\ref{prop:donor-monotonicity}. It enlarges the exact robust safe set
only by changing the observation or compatible-world class so that a previously
unsafe history enters $\mathcal H_{\rm safe}$. For the operational rule $C^+$,
this can occur by splitting a mixed closure fibre, discharging an obligation,
establishing a missing validity condition $V$, or making a required action and
terminal reachable under the registered budget/interface. The list is
exhaustive only relative to those declared state coordinates.
\end{corollary}

\begin{corollary}[Comparator-adapter admissibility]
Let $T$ be a comparator's complete native transcript and let $a(T)$ be the
normalized transcript supplied to the frozen evaluation. An adapter is lawful
only if $a$ is a prospectively fixed deterministic map. Thus under every
declared world law $W\to T\to a(T)$ is a Markov chain, with no new query,
source, world signal or unregistered action introduced by $a$. It must map
native actions, costs, outputs and terminals into the frozen
route/action/output namespace while preserving every registered terminal
distinction. Other than prospectively fixed identity normalization, it cannot
invent source identities or content; in particular, a native empty result,
error or timeout remains non-evidence-bearing and cannot become acquired
evidence, an obligation discharge, or task closure. Under these conditions
every decision based on $a(T)$ also factors through $T$, so the adapter cannot
split a native transcript fibre or lift the acquisition ceiling. It may only
preserve or coarsen the comparator's information.

Adapter lawfulness does not itself authorize closure. Soundness of a close
terminal still requires the sufficient condition
$V(h)\wedge R(h)=\varnothing\Rightarrow h\in\mathcal H_{\rm safe}$, and equality
with maximal robust authority still requires the reverse implication. Thus a
native answer or end terminal may be recorded faithfully without being
promoted to a P2 task-level \textsc{close} terminal.
\end{corollary}

This is the sense in which the proposed theory envelops its nearest work.
Retrieval, coverage, information utility, stopping, obligation certificates and
memory are all admissible mechanisms inside the envelope.  Their scientific
value is measured by the joint frontier they make reachable; their outputs do
not acquire global authority merely by being numerically confident.

\subsection{Safe residual envelopes and the donor-saturation boundary}

Fix a source-bound finite family $F$, exact donor $d$, residual $r$ and a
predeclared conjunction $\mathcal K(F)$ of CRE20, R@10, WSS95, sign, harm and
absolute-work-saving gates. Let $B_F$ collect source, population,
implementation and custody bindings, with every unbound predicate treated as
false. The certificate-switched policy is
\[
 \pi_r^{\mathrm{env}}=
 \begin{cases}
 e_r,&B_F=1\ \text{and}\ x(r)\in\mathcal K(F),\\
 d,&\text{otherwise}.
 \end{cases}
\]
Thus an unavailable or rejected residual is not imputed as a tie or zero; it
executes the exact donor.

\begin{theorem}[Simultaneously certified residual embedding]
\label{thm:safe-residual-envelope}
Suppose the family, donor, residual, endpoints and gates are fixed before
candidate outcomes. If a simultaneous $(1-\alpha)$ confidence set for the
complete gate vector is contained in $\mathcal K(F)$, then the probability of
admitting a target vector outside $\mathcal K(F)$ is at most $\alpha$.
Moreover, the attainable set of the certificate-switched class contains the
donor attainable set, and every rejected or unbound residual executes $d$.
\end{theorem}

\begin{proof}
False admission requires the true vector to lie outside $\mathcal K(F)$ while
the confidence set lies inside it, which is possible only on simultaneous
noncoverage. Donor containment follows by construction, and fallback is the
second branch of the switch.
\end{proof}

The conjunction is noncompensatory: a sufficiently large gain on one endpoint
cannot buy a failed coprimary, sign or harm gate. In particular, V10's positive
mean WSS95 cannot repair its failed CRE20/R@10 magnitude and sign gates.

Now let $s_D(w)$ be all information and retained state available to a donor
class $\Pi_D$ in world $w$, and let $\mathcal A_D(w)$ be its admissible action
and acquisition support under the matched contract. Call $\Pi_D$ saturated if
it contains every allowed policy measurable from $s_D$ that uses only
$\mathcal A_D$.

\begin{theorem}[Donor-fibre saturation boundary]
\label{thm:donor-fibre-saturation}
Under the same task, access, budget and authority contract, a candidate whose
action is constant on every fibre of $s_D$ and lies in $\mathcal A_D$ belongs
to a saturated $\Pi_D$. Therefore strict outside-envelope separation requires
at least one of: a signal or retained state that separates a donor fibre, an
acquisition action outside donor support, or proof that the declared class was
not saturated. Changing the cost, access, world or authority contract changes
the estimand and does not establish dominance under the original contract.
\end{theorem}

\begin{proof}
Fibre constancy gives a well-defined factorization $c=g\circ s_D$. Matched
action support and contract guards make this an admissible donor policy, which
saturation includes. The separation statement is the contrapositive; a contract
mismatch instead violates the comparison premise.
\end{proof}

Finite residual failures do not establish saturation or universal donor
optimality: absent a cross-world restriction, two extensions can agree on all
observed reviews while a new-family residual fails every gate in one and passes
them in the other. The title-emphasis transform is deterministic from fields
already in donor state, so it adds no information; V10 rejects that one
algorithmic transform, not every possible same-information learner.

\subsection{A failure-derived candidate: obligation-directed adaptive
exploration}

The theory also specifies a materially stronger acquisition candidate rather
than treating fail-closed authority as a substitute for finding evidence.  Let
$m_t(o)$ be the posterior unresolved mass of material obligation $o$, let
$Q_t(a)$ be the prospectively evaluated option value that action $a$ preserves
for hidden future questions, and let $J_t(a)$ be the expected reduction in
closure-fibre ambiguity---the probability mass of closable/open worlds that the
action is expected to separate.  An \emph{obligation-directed adaptive
exploration} (ODAE) policy chooses from route invention, retrieval, processing,
contradiction resolution, state preservation/recovery and closure verification
by a charged index of the form
\[
 a_t\in\arg\max_{a\in\mathcal A(h_t)}
 \frac{
   \operatorname{LCB}_t\!\left[
     \Delta\sum_{o\in R(h_t)}\mu_o m_t(o)\mid a
   \right]
   +\lambda Q_t(a)+\beta J_t(a)
 }{\operatorname{cost}_t(a)}.
\]
The lower-confidence term prevents an uncertain route estimate from being
treated as established gain; $Q_t$ prevents current-question compression from
destroying future route cues; and $J_t$ values actions that make the closure
decision identifiable even when their immediate document count is low.  If no
action has validated support through the frozen interface, the controller emits
\textsc{cannot-check} rather than assigning that route zero value.

Generic value-of-information, active search, coverage estimation and adaptive
stopping are donor-owned.  The proposed residual is their non-compensatory
coupling to material scientific obligations, option preservation and closure
fibre separation under one charged budget.  ODAE is a prospective policy family,
not an executed result or current general capability.  Its direct falsifier is a
donor-complete controller that reaches the same or a better guarded frontier on
fresh tasks.  Its causal tests must separately vary (i) route support, (ii)
retrieved-but-unprocessed evidence, (iii) contradiction sides, (iv) future-query
optionality and (v) closable clean cases; improvement on only the family that
inspired the policy is insufficient.

\subsection{A prospective cross-arena discriminator}

The theory converts the historical adverse terminals into explicit research
problems rather than deleting them.  The 7/229 exposure result demands route
invention and admissible-support estimation before reranking.  The one-third
scoring-eligible exposure mismatch demands source-interface matching rather
than call-count matching.  The retrieved-but-unprocessed and future-query
failures demand option-valued state preservation.  Provider-invalid and
unresolved-route failures demand independent observability custody and an
authority-complete state.  Clean closable tasks remain a required counterweight
against an always-open policy.

A separately frozen successor gives these problems a common empirical target.
It specifies 768 source-disjoint task-world clusters---192 in each of four
strata---across four arenas, identical admissible route exposure, three frozen
primary comparators, an intersection requirement over every contrast, and a
joint power target of 0.90 under its declared planning model.  That design is a
protocol, not a result.  Archived baseline-task-world reproduction, four
source-disjoint arena snapshots, a capability-matched external comparator,
independent evaluator/gold custody, and live response/cost receipts are still
required.  Until those bindings exist and the protected campaign runs, the
general cross-arena envelope-superiority hypothesis remains untested and the
historical manuscript terminal remains unchanged.
